\documentclass[10pt,a4paper]{article}
\usepackage[utf8]{inputenc}
\usepackage[a4paper,left=3cm,right=2cm]{geometry}
\usepackage{amsmath}
\usepackage[thmmarks, amsmath, thref]{ntheorem}
\usepackage{amsfonts}
\usepackage{amssymb}
\usepackage{fancyhdr}
\pagestyle{fancy}
\usepackage{anysize}
\usepackage{graphicx}
\usepackage{float}
\usepackage{hyperref}
\usepackage{multirow}
\usepackage{enumitem}
\newtheorem{theorem}{Theorem}
\theoremsymbol{\ensuremath{\blacksquare}}
\newtheorem*{solution}{Solución:}

\usepackage{xcolor}

% Margenes y formalidades

\textwidth 6.25in % Margen Derecho y ancho del texto
\textheight 8.5in
\oddsidemargin 0in % Margen Izquierdo
\headheight 0in % Largo del texto
\leftmargin 1in
\parindent 0em % Salto entre la indentación (Sangría)https://www.overleaf.com/4519127669sdpptkhvsbfn
\topmargin -0.5in % Margen Superior
\parskip 0.5ex % Salto entre parrafos
\baselineskip 1pt

%\lhead{\includegraphics[scale=0.2]{logo.png}} % texto izquierda de la cabecera
%\chead{}                                      % texto centro de la cabecera
%\rhead{\includegraphics[scale=0.2]{FIG/logo-mba-2016.png}} % número de página a la derecha

\renewcommand{\headrulewidth}{0.4pt} % grosor de la línea de la cabecera
\renewcommand{\footrulewidth}{0.4pt} % grosor de la línea del pie

\begin{document}



\pagebreak
\vspace*{4mm} 

\begin{center}
\begin{huge}
\textbf{\\Ejercicios Econometría 2}
\end{Large}\\
%{\large 2023-2}\\
\end{center}


\section{Muestreo e Intervalos de Confianza}

Un investigador quiere estimar el gasto promedio mensual en transporte de los estudiantes de una universidad. Para ello, selecciona una muestra aleatoria de 100 estudiantes y obtiene los siguientes datos:

\begin{itemize}
    \item \textbf{Media poblacional} ($\mu$) = 120 dólares (dado como referencia)
    \item \textbf{Varianza poblacional} ($\sigma^2$) = 900 dólares
    \item \textbf{Media muestral} ($\bar{x}$) = 125 dólares
    \item \textbf{Tamaño de la muestra} ($n$) = 100
    \item \textbf{Nivel de confianza} = 95\%
\end{itemize}

\textbf{Preguntas:}

\begin{enumerate}
    \item Calcula el intervalo de confianza del 95\% para la media poblacional.
    \item Interpreta el intervalo de confianza en el contexto del problema.
    \item Determina el tamaño mínimo de la muestra ($n$) necesario para que la media poblacional de 120 no pertenezca al intervalo de confianza del 95\% basado en la media muestral de 125.
\end{enumerate}

\section{Regresión Lineal Simple con suma de Constante}

Sean \(\hat{\beta}_0\) y \(\hat{\beta}_1\) el intercepto y la pendiente en la regresión de \(y_i\) sobre \(x_i\), usando \(n\) observaciones.

Ahora, sean \(\tilde{\beta}_0\) y \(\tilde{\beta}_1\) la intersección con el eje \(y\) y la pendiente en la regresión de \((c_1 + y_i)\) sobre \((c_2 + x_i)\) (sin ninguna restricción sobre \(c_1\) o \(c_2\)). Muestre que \(\tilde{\beta}_1 = \hat{\beta}_1\) y \(\tilde{\beta}_0 = \hat{\beta}_0 +c_1- c_2 \cdot \hat{\beta}_1\). Compare los $R^2$ de ambas regresiones.




\section{Regresión Lineal Simple: Sesgo por error de medición}

Sean \( \hat{\beta}_0 \) y \( \hat{\beta}_1 \) los parámetros estimados de la regresión de \( y_i \) sobre \( x_i \), y sean \( \tilde{\beta}_0 \) y \( \tilde{\beta}_1 \) los parámetros estimados de la regresión de \( y_i \) sobre \( (u + x_i) \), donde \( u \sim N(0, 1) \). 
\begin{enumerate}
    \item Encuentre los nuevos valores de \( \tilde{\beta}_0 \) y \( \tilde{\beta}_1 \) en función de \( \hat{\beta}_0 \) y \( \hat{\beta}_1 \).

    \item Supongamos que tenemos una muestra de 1000 datos con $\Bar{x} = 5$, $\Bar{y}=10$, $Var(x) = 4$ y $Cov(x,y)=3$, encuentre el valor de los estimadores \( \hat{\beta}_0 \) y \( \hat{\beta}_1 \) y de \( \tilde{\beta}_0 \) y \( \tilde{\beta}_1 \), y las ecuaciones que describen al regresando.

\end{enumerate}



\section{Sesgo de selección}
Supongamos que un gobierno implementa un programa de capacitación laboral para ayudar a las personas desempleadas a encontrar trabajo. El programa ofrece cursos gratuitos de habilidades técnicas y entrevistas simuladas. Sin embargo, la participación en el programa es voluntaria, y las personas deciden inscribirse o no en función de sus propias características (por ejemplo, motivación, nivel educativo previo, etc.).

Tenemos los siguientes datos para dos grupos: \textbf{tratados} (participantes del programa) y \textbf{no tratados} (no participantes).

\begin{center}
\begin{tabular}{|c|c|c|c|}
\hline
Individuo & Tratamiento (\( D_i \)) & Salario (\( Y_i \)) & Nivel educativo (\( X_i \)) \\
\hline
1 & 1 & 2500 & 12 \\
2 & 1 & 3000 & 14 \\
3 & 1 & 2800 & 13 \\
4 & 0 & 2000 & 10 \\
5 & 0 & 2200 & 11 \\
6 & 0 & 2100 & 10 \\
\hline
\end{tabular}
\end{center}

\begin{enumerate}
    \item Comente cuales serían los posibles sesgos que existen en comparar el grupo de tratados vs los no tratados.

    \item Calcule el $\beta$ que asimila el ATE y ATT. 

    \item ¿Existe un sesgo dado el nivel educativo?

    \item supongamos que incluimos la variable nivel educativo en una regresión lineal dandonos la siguiente ecuación: 
    $$\hat{Y}_i = 1500+200D_i + 100X_i$$
    ¿Cúal es el valor del sesgo de selección? ¿Cúal sería el impacto real de la participación en el programa en el salario?
\end{enumerate}



\section{Sesgo por variables omitidas}
Queremos estudiar el efecto de la educación (\( X_1 \)) sobre el salario (\( Y \)). Sin embargo, omitimos la variable "habilidades innatas" (\( X_2 \)), que también afecta el salario y está correlacionada con la educación. Esto introduce un \textbf{sesgo por variables omitidas}.

\textbf{Datos}
Tenemos los siguientes datos para 6 individuos:

\begin{center}
\begin{tabular}{|c|c|c|c|}
\hline
Individuo & Educación (\( X_1 \)) & Habilidades innatas (\( X_2 \)) & Salario (\( Y \)) \\
\hline
1 & 12 & 5 & 2500 \\
2 & 14 & 7 & 3000 \\
3 & 16 & 6 & 3500 \\
4 & 10 & 4 & 2000 \\
5 & 13 & 5 & 2800 \\
6 & 15 & 8 & 4000 \\
\hline
\end{tabular}
\end{center}

\begin{enumerate}
    \item Estimando el efecto de la educación (\( X_1 \)) sobre el salario (\( Y \)) omitiendo la variable "habilidades innatas" (\( X_2 \)) (regresión lineal simple), tenemos los siguientes resultados:
    \begin{itemize}
    \item Coeficiente de educación (\( \hat{\beta}_1 \)): 200.
    \item Intercepto (\( \hat{\beta}_0 \)): 500.
    \item \( R^2 \): 0.75.
    \item \( SSR \): 150,000.
    \end{itemize}
    Luego, incluyendo la variable "habilidades innatas" (\( X_2 \)) en el modelo (regresión múltiple), obtenemos los siguientes resultados:
    \begin{itemize}
    \item Coeficiente de educación (\( \hat{\beta}_1 \)): 150.
    \item Coeficiente de habilidades innatas (\( \hat{\beta}_2 \)): 100.
    \item Intercepto (\( \hat{\beta}_0 \)): 300.
    \item \( R^2 \): 0.90.
    \item \( SSR \): 50,000.
\end{itemize}
Compare los resultados de ambos modelos e identifique el signo del sesgo.

    \item Explica por qué ocurre el sesgo por variables omitidas.
\end{enumerate}



    




\end{document}